% id: 2022-ualype-1-1 | fonte: listas/2022/Ualype/lista2022ualype1.pdf, p. 1
\textbf{(a)} Demonstre a equação de continuidade na forma mostrada, e então prove que $\nabla\cdot\vec{v} = 0$ para um escoamento de um fluido incompressível. Use o fato de que, dentro de um volume de controle fixo, massa não pode ser criada nem destruída. Utilize o Teorema de Gauss: $\iiint_V (\nabla\cdot\mathbf{F})dV = \oint_S \mathbf{F}\cdot d\vec{S}$

\textbf{(b)} Demonstre a Equação de Bernoulli, válida ao longo de qualquer linha de corrente em um escoamento estacionário. Use o fato de que, dentro de um volume de controle, energia não pode ser criada nem destruída.

\textbf{(c)} Fora do regime estacionário, é necessária uma correção ao Teorema de Bernoulli. Assumindo um escoamento irrotacional ($\nabla\times\vec{v} = 0$), é possível atribuir ao campo de velocidades um potencial $\phi$ tal que $\vec{v} = -\nabla\phi$. Com isso, mostre que
\[ \frac{p}{\rho} + \Psi + \frac{v^2}{2} - \frac{\partial\phi}{\partial t} = f(t) \]
